Геологические среды обладают выраженной неоднородностью: слоистостью, анизотропией, вариацией температуры и плотности. В случае пространственных вариаций параметров Ламе уравнение движения  принимает вид:

\begin{equation}
\rho(x) \frac{\partial^2 u_i}{\partial t^2} =
\frac{\partial}{\partial x_j}
\left[
\lambda(x) \delta_{ij} \frac{\partial u_k}{\partial x_k}
+ \mu(x) \left(
\frac{\partial u_i}{\partial x_j} + \frac{\partial u_j}{\partial x_i}
\right)
\right]
- \gamma(x) \frac{\partial T}{\partial x_i}.
\end{equation}

Анизотропия может быть учтена через тензор жесткости $C_{ijkl}$:

\begin{equation}
\sigma_{ij} = C_{ijkl} \epsilon_{kl} - \gamma_{ij} T.
\end{equation}

Моделирование таких сред требует высокопроизводительных численных методов: FDTD, спектральных методов или гибридных схем.
